% Figure 3.1, the dynamics of one instance over three decision epochs.
% Schematic. The geometry and the commitments are invented for clarity and are
% not drawn from any run.
% R2 is the thread the reader follows. R2, the two tasks it takes and its
% busy-time bar are drawn at full weight in the accent colour, and everything
% else is present but held back in a light grey at reduced stroke weight.
% Two rules govern the marks. A task carries its state on its outline alone,
% one of three, and a dash means one thing only, being the carry leg a robot is
% on between a pickup and its drop-off. A delivered task is removed from the
% panel rather than given a state of its own, which clears the later panels.
% Preamble requirements
%   \usepackage{tikz}
%   \usetikzlibrary{arrows}
% Natural size is 14.48 cm by 7.90 cm, sized for the thesis text width.
% Include with % Figure 3.1, the dynamics of one instance over three decision epochs.
% Schematic. The geometry and the commitments are invented for clarity and are
% not drawn from any run.
% R2 is the thread the reader follows. R2, the two tasks it takes and its
% busy-time bar are drawn at full weight in the accent colour, and everything
% else is present but held back in a light grey at reduced stroke weight.
% Two rules govern the marks. A task carries its state on its outline alone,
% one of three, and a dash means one thing only, being the carry leg a robot is
% on between a pickup and its drop-off. A delivered task is removed from the
% panel rather than given a state of its own, which clears the later panels.
% Preamble requirements
%   \usepackage{tikz}
%   \usetikzlibrary{arrows}
% Natural size is 14.48 cm by 7.90 cm, sized for the thesis text width.
% Include with % Figure 3.1, the dynamics of one instance over three decision epochs.
% Schematic. The geometry and the commitments are invented for clarity and are
% not drawn from any run.
% R2 is the thread the reader follows. R2, the two tasks it takes and its
% busy-time bar are drawn at full weight in the accent colour, and everything
% else is present but held back in a light grey at reduced stroke weight.
% Two rules govern the marks. A task carries its state on its outline alone,
% one of three, and a dash means one thing only, being the carry leg a robot is
% on between a pickup and its drop-off. A delivered task is removed from the
% panel rather than given a state of its own, which clears the later panels.
% Preamble requirements
%   \usepackage{tikz}
%   \usetikzlibrary{arrows}
% Natural size is 14.48 cm by 7.90 cm, sized for the thesis text width.
% Include with % Figure 3.1, the dynamics of one instance over three decision epochs.
% Schematic. The geometry and the commitments are invented for clarity and are
% not drawn from any run.
% R2 is the thread the reader follows. R2, the two tasks it takes and its
% busy-time bar are drawn at full weight in the accent colour, and everything
% else is present but held back in a light grey at reduced stroke weight.
% Two rules govern the marks. A task carries its state on its outline alone,
% one of three, and a dash means one thing only, being the carry leg a robot is
% on between a pickup and its drop-off. A delivered task is removed from the
% panel rather than given a state of its own, which clears the later panels.
% Preamble requirements
%   \usepackage{tikz}
%   \usetikzlibrary{arrows}
% Natural size is 14.48 cm by 7.90 cm, sized for the thesis text width.
% Include with \input{figure_worked_instance} and do not rescale.
%
\providecommand{\figlab}{\fontsize{8}{9.4}\selectfont\hyphenpenalty=10000\relax}%
\providecommand{\figann}{\fontsize{7}{8.2}\selectfont}%
%
\providecolor{fgAccent}{RGB}{20,86,150}%
\providecolor{fgAccentFill}{RGB}{224,235,246}%
\providecolor{fgDark}{RGB}{45,45,45}%
\providecolor{fgMid}{RGB}{120,120,120}%
\providecolor{fgLine}{RGB}{160,160,160}%
\providecolor{fgFill}{RGB}{243,243,243}%
\providecolor{fgBand}{RGB}{247,247,247}%
%
\begin{tikzpicture}[
  x=1cm, y=1cm,
  >=stealth',
  line join=round,
  every node/.style={font=\figann},
  % arena frame
  arena/.style={draw=fgLine, line width=0.5pt, fill=fgBand, rounded corners=1.5pt},
  %
  % ---------------------------------------------------------------- emphasis
  % Two families of the same shapes. The f family is R2's thread, in the accent
  % colour at full weight. The b family is everything else, in a light grey at
  % reduced weight. Shape, fill and outline carry the same meaning in both, so
  % the emphasis never changes what a marker means.
  %
  % robots. An open square is available to the dispatcher this epoch, a filled
  % square is working and cannot be assigned.
  fRfree/.style={rectangle, draw=fgAccent, line width=1.0pt, fill=white,
    minimum size=2.4mm, inner sep=0pt},
  fRbusy/.style={rectangle, draw=fgAccent, line width=0.7pt, fill=fgAccent,
    minimum size=2.4mm, inner sep=0pt},
  bRfree/.style={rectangle, draw=fgLine, line width=0.5pt, fill=white,
    minimum size=2.1mm, inner sep=0pt},
  bRbusy/.style={rectangle, draw=fgLine, line width=0.5pt, fill=fgLine,
    minimum size=2.1mm, inner sep=0pt},
  %
  % tasks. Circular throughout. The three states are set apart by rendering
  % rather than by stroke weight alone, so they stay separable at 1.6 mm.
  %   not yet released  one ring,  drawn at reduced opacity, ghosted
  %   pending           one ring,  full weight
  %   being carried     two rings, full weight, largest
  % Fill carries pickup against drop-off, tinted against open. The tint is kept
  % well below the outline colour so the outline stays readable on a pickup.
  % Nothing here is dashed.
  fRelP/.style={circle, draw=fgAccent, line width=1.0pt, fill=fgAccent!45,
    minimum size=1.6mm, inner sep=0pt, opacity=0.45},
  fRelD/.style={circle, draw=fgAccent, line width=1.0pt, fill=white,
    minimum size=1.6mm, inner sep=0pt, opacity=0.45},
  fPendP/.style={circle, draw=fgAccent, line width=1.0pt, fill=fgAccent!45,
    minimum size=1.7mm, inner sep=0pt},
  fPendD/.style={circle, draw=fgAccent, line width=1.0pt, fill=white,
    minimum size=1.7mm, inner sep=0pt},
  fCarryP/.style={circle, draw=fgAccent, line width=0.75pt, double=white,
    double distance=0.5pt, fill=fgAccent!45, minimum size=2.1mm, inner sep=0pt},
  fCarryD/.style={circle, draw=fgAccent, line width=0.75pt, double=white,
    double distance=0.5pt, fill=white, minimum size=2.1mm, inner sep=0pt},
  bRelP/.style={circle, draw=fgLine, line width=0.8pt, fill=fgLine!35,
    minimum size=1.5mm, inner sep=0pt, opacity=0.45},
  bRelD/.style={circle, draw=fgLine, line width=0.8pt, fill=white,
    minimum size=1.5mm, inner sep=0pt, opacity=0.45},
  bPendP/.style={circle, draw=fgLine, line width=0.8pt, fill=fgLine!35,
    minimum size=1.6mm, inner sep=0pt},
  bPendD/.style={circle, draw=fgLine, line width=0.8pt, fill=white,
    minimum size=1.6mm, inner sep=0pt},
  bCarryP/.style={circle, draw=fgLine, line width=0.5pt, double=white,
    double distance=0.4pt, fill=fgLine!35, minimum size=1.9mm, inner sep=0pt},
  bCarryD/.style={circle, draw=fgLine, line width=0.5pt, double=white,
    double distance=0.4pt, fill=white, minimum size=1.9mm, inner sep=0pt},
  %
  % the solid line pairing a pickup with its drop-off, ghosted to match while
  % the task it belongs to has not yet been released
  flink/.style={draw=fgAccent, line width=0.5pt},
  blink/.style={draw=fgLine, line width=0.3pt},
  fghost/.style={draw=fgAccent, line width=0.5pt, opacity=0.45},
  bghost/.style={draw=fgLine, line width=0.3pt, opacity=0.45},
  % the carry leg, the only dashed mark in the figure. It replaces the pairing
  % line for as long as a robot is on it.
  fcarry/.style={draw=fgAccent, line width=0.7pt,
    dash pattern=on 2.4pt off 1.5pt},
  bcarry/.style={draw=fgLine, line width=0.4pt,
    dash pattern=on 2.4pt off 1.5pt},
  % a commitment made this epoch, robot to pickup
  commit/.style={draw=fgAccent, line width=1.0pt},
  bcommit/.style={draw=fgLine, line width=0.5pt},
  % the leaders that name a pickup, a drop-off and an available robot once, in
  % the first panel, in place of legend entries for them
  leader/.style={draw=fgMid, line width=0.3pt},
  % bars
  fbar/.style={draw=fgAccent, line width=0.8pt, fill=fgAccent},
  bbar/.style={draw=fgLine, line width=0.4pt, fill=fgFill},
  guide/.style={draw=fgLine, line width=0.4pt,
    dash pattern=on 0.1pt off 1.8pt, line cap=round},
  span/.style={draw=fgAccent, line width=0.9pt},
  % text
  fTag/.style={text=fgAccent, inner sep=1.5pt, font=\figann},
  bTag/.style={text=fgMid, inner sep=1.5pt, font=\figann},
  % the single annotation set in black. It marks the point of the middle panel
  % without joining R2's accent thread, which stays reserved for R2.
  dTag/.style={text=fgDark, inner sep=1.5pt, font=\figann},
  fLab/.style={text=fgAccent, inner sep=1.5pt, font=\figann},
  bLab/.style={text=fgMid, inner sep=1.5pt, font=\figann},
  head/.style={text=fgDark, inner sep=0pt, font=\figlab}
]

\path[use as bounding box] (0,-7.55) rectangle (14.48,0.35);

% =================================================================== panels
% Panel origins. Three arenas of 4.50 cm separated by 0.49 cm, so the row is
% exactly the 14.48 cm text width. Every panel uses the same local grid of
% 10 by 11 units at 0.45 cm per unit, so a position means the same thing in all
% three and the robots can be tracked across the row.
% Four tasks. R1 takes task 1 at the first epoch, R2 takes task 2 at the first
% epoch and task 3 at the third, and R3 takes task 4 at the third. At the
% second epoch R3 is free and task 3 is pending, and R3 is left uncommitted
% anyway, which is a move a greedy assignment cannot make. R2's own position
% carries the episode.
% It starts away from the work, then sits on task 2's carry leg while it
% executes, then stands at task 2's drop-off, free again. Task 2 is delivered
% by the third panel and so is no longer drawn there.

% ------------------------------------------------------------------ epoch 1
\begin{scope}[shift={(0.00,0)}]
  \node[anchor=south west, head] at (0,0.06) {Epoch 1};
  \path[arena] (0,-4.95) rectangle (4.50,0);
  \begin{scope}[shift={(0,-4.95)}, scale=0.45]
    % pairing lines
    \draw[blink] (3.7,9.5) -- (8.7,8.4);
    \draw[flink] (2.4,4.8) -- (4.8,5.9);
    \draw[bghost] (6.6,1.6) -- (9.1,3.2);
    \draw[fghost] (6.3,6.8) -- (8.6,7.3);
    % task 1, pending, R1's
    \node[bPendP] at (3.7,9.5) {}; \node[bPendD] at (8.7,8.4) {};
    \node[bTag, anchor=south] at (3.7,9.72) {1};
    % task 2, pending, R2's
    \node[fPendP] at (2.4,4.8) {}; \node[fPendD] at (4.8,5.9) {};
    \node[fTag, anchor=east] at (2.12,4.8) {2};
    % task 4, not yet released, R3's later
    \node[bRelP] at (6.6,1.6) {}; \node[bRelD] at (9.1,3.2) {};
    \node[bTag, anchor=north] at (6.6,1.36) {4};
    % task 3, not yet released, R2's later
    \node[fRelP] at (6.3,6.8) {}; \node[fRelD] at (8.6,7.3) {};
    \node[fTag, anchor=south] at (6.3,7.04) {3};
    % commitments
    \draw[bcommit,->] (1.856,9.008) -- (3.371,9.412);
    \draw[commit,->] (1.703,3.772) -- (2.220,4.535);
    % the three marks that are named here rather than in the legend
    \draw[leader] (2.4,5.60) -- (2.4,5.02);
    \node[bTag, anchor=south] at (2.4,5.60) {pickup};
    \draw[leader] (4.8,5.30) -- (4.8,5.71);
    \node[bTag, anchor=north] at (4.8,5.30) {drop-off};
    \draw[leader] (3.6,1.05) -- (3.6,1.37);
    \node[bTag, anchor=north] at (3.6,1.05) {available};
    % robots, all three available and none of them has begun. R1 and R2 are set
    % in far enough that their labels clear the left edge of the bounding box.
    \node[bRfree] at (1.45,8.9) {}; \node[bLab, anchor=east] at (1.15,8.9) {R1};
    \node[fRfree] at (1.45,3.4) {}; \node[fLab, anchor=east] at (1.13,3.4) {R2};
    \node[bRfree] at (3.6,1.6) {}; \node[bLab, anchor=east] at (3.30,1.6) {R3};
  \end{scope}
\end{scope}

% ------------------------------------------------------------------ epoch 2
\begin{scope}[shift={(4.99,0)}]
  \node[anchor=south west, head] at (0,0.06) {Epoch 2};
  \path[arena] (0,-4.95) rectangle (4.50,0);
  \begin{scope}[shift={(0,-4.95)}, scale=0.45]
    % carry legs, the only dashed marks. R1 and R2 are on theirs.
    \draw[bcarry] (3.7,9.5) -- (8.7,8.4);
    \draw[fcarry] (2.4,4.8) -- (4.8,5.9);
    % pairing lines
    \draw[blink] (6.6,1.6) -- (9.1,3.2);
    \draw[flink] (6.3,6.8) -- (8.6,7.3);
    % task 1, being carried
    \node[bCarryP] at (3.7,9.5) {}; \node[bCarryD] at (8.7,8.4) {};
    % task 2, being carried, R2's
    \node[fCarryP] at (2.4,4.8) {}; \node[fCarryD] at (4.8,5.9) {};
    \node[fTag, anchor=east] at (2.12,4.8) {2};
    % task 4, released this epoch and left waiting
    \node[bPendP] at (6.6,1.6) {}; \node[bPendD] at (9.1,3.2) {};
    % task 3, released this epoch and left waiting, R2's later
    \node[fPendP] at (6.3,6.8) {}; \node[fPendD] at (8.6,7.3) {};
    \node[fTag, anchor=south] at (6.3,7.04) {3};
    \node[fTag, anchor=north] at (6.3,6.56) {waits};
    % no commitment is made this epoch
    % robots. R1 and R2 sit on their carry legs, so they are mid-execution and
    % cannot be assigned. R3 is at its start position and available, and is left
    % uncommitted even though task 3 is pending. Nothing here claims that
    % holding R3 back is optimal or that it is what any trained policy does.
    % It is drawn only to show that the decision is available to a dispatcher
    % and is not available to a greedy assignment.
    \node[bRbusy] at (5.7,9.06) {}; \node[bLab, anchor=south] at (5.7,9.36) {R1};
    \node[fRbusy] at (3.48,5.295) {};
    \node[fLab, anchor=north] at (3.48,5.02) {R2};
    \node[bRfree] at (3.6,1.6) {}; \node[bLab, anchor=east] at (3.30,1.6) {R3};
    \draw[leader] (3.6,1.05) -- (3.6,1.37);
    \node[dTag, anchor=north] at (3.6,1.05) {held back};
  \end{scope}
\end{scope}

% ------------------------------------------------------------------ epoch 3
\begin{scope}[shift={(9.98,0)}]
  \node[anchor=south west, head] at (0,0.06) {Epoch 3};
  \path[arena] (0,-4.95) rectangle (4.50,0);
  \begin{scope}[shift={(0,-4.95)}, scale=0.45]
    % carry legs
    \draw[bcarry] (3.7,9.5) -- (8.7,8.4);
    % pairing lines
    \draw[blink] (6.6,1.6) -- (9.1,3.2);
    \draw[flink] (6.3,6.8) -- (8.6,7.3);
    % task 1, still being carried
    \node[bCarryP] at (3.7,9.5) {}; \node[bCarryD] at (8.7,8.4) {};
    % task 2 has been delivered and is no longer drawn
    % task 4, the task R3 was held back from, taken this epoch by R3
    \node[bPendP] at (6.6,1.6) {}; \node[bPendD] at (9.1,3.2) {};
    % task 3, the waiting task, taken this epoch by R2
    \node[fPendP] at (6.3,6.8) {}; \node[fPendD] at (8.6,7.3) {};
    \node[fTag, anchor=south] at (6.3,7.04) {3};
    % commitments. R2 has just delivered, and R3 was free through the last
    % epoch and takes the task it was held back from.
    \draw[commit,->] (5.170,6.122) -- (6.035,6.641);
    \draw[bcommit,->] (4.02,1.6) -- (6.26,1.6);
    % robots. R1 is on its carry leg. R2 stands at task 2's drop-off, which is
    % where delivering it left it, and is free again. R3 has not moved.
    \node[bRbusy] at (7.7,8.62) {}; \node[bLab, anchor=south] at (7.7,8.92) {R1};
    \node[fRfree] at (4.8,5.9) {}; \node[fLab, anchor=east] at (4.48,5.9) {R2};
    \node[bRfree] at (3.6,1.6) {}; \node[bLab, anchor=east] at (3.30,1.6) {R3};
  \end{scope}
\end{scope}

% ===================================================================== bars
% Cumulative busy time after the three epochs. R2 has finished one task and is
% starting a second, so it has done the most. R1 is still on the single long
% task it took at the first epoch, and R3 stood idle through the first two
% epochs and has done the least. The bars were scaled down by a common factor
% to clear the legend column on the right, so their relative lengths are
% unchanged. The scale is 1.03 cm per unit of busy time.
\node[anchor=south west, head] at (0,-5.40) {Cumulative busy time};

\path[bbar] (0.85,-5.94) rectangle (8.53,-5.70);
\node[bLab, anchor=east] at (0.75,-5.82) {R1};

\path[fbar] (0.85,-6.36) rectangle (10.15,-6.12);
\node[fLab, anchor=east] at (0.75,-6.24) {R2};

\path[bbar] (0.85,-6.78) rectangle (5.75,-6.54);
\node[bLab, anchor=east] at (0.75,-6.66) {R3};

% the balance term, read off as the distance between the longest and the
% shortest bar. R2 is the longest, so it is one end of the span.
\draw[guide] (10.15,-6.36) -- (10.15,-7.14);
\draw[guide] (5.75,-6.78) -- (5.75,-7.14);
\draw[span,<->] (5.75,-7.14) -- (10.15,-7.14);
\node[anchor=north, inner sep=1.5pt, text=fgAccent] at (7.95,-7.20) {balance};

% =================================================================== legend
% A column of three set beside the bars rather than under them, centred on the
% middle bar so the two blocks read as one band. The three task states and
% nothing else. Pickup, drop-off and an available robot are named directly in
% the first panel instead. The swatches use the accent family, which is the
% unfaded rendering of every mark in the figure. Only the being-carried swatch
% is dashed. The column starts at 11.30, clear of the longest bar at 10.15 and
% of the balance span and its label, which sit lower and further left.
\draw[fghost] (11.30,-5.74) -- (11.90,-5.74);
\node[fRelP] at (11.30,-5.74) {}; \node[fRelD] at (11.90,-5.74) {};
\node[anchor=west, bTag] at (12.05,-5.74) {not yet released};

\draw[flink] (11.30,-6.24) -- (11.90,-6.24);
\node[fPendP] at (11.30,-6.24) {}; \node[fPendD] at (11.90,-6.24) {};
\node[anchor=west, bTag] at (12.05,-6.24) {pending};

\draw[fcarry] (11.30,-6.74) -- (11.90,-6.74);
\node[fCarryP] at (11.30,-6.74) {}; \node[fCarryD] at (11.90,-6.74) {};
\node[anchor=west, bTag] at (12.05,-6.74) {being carried};

\end{tikzpicture}%
 and do not rescale.
%
\providecommand{\figlab}{\fontsize{8}{9.4}\selectfont\hyphenpenalty=10000\relax}%
\providecommand{\figann}{\fontsize{7}{8.2}\selectfont}%
%
\providecolor{fgAccent}{RGB}{20,86,150}%
\providecolor{fgAccentFill}{RGB}{224,235,246}%
\providecolor{fgDark}{RGB}{45,45,45}%
\providecolor{fgMid}{RGB}{120,120,120}%
\providecolor{fgLine}{RGB}{160,160,160}%
\providecolor{fgFill}{RGB}{243,243,243}%
\providecolor{fgBand}{RGB}{247,247,247}%
%
\begin{tikzpicture}[
  x=1cm, y=1cm,
  >=stealth',
  line join=round,
  every node/.style={font=\figann},
  % arena frame
  arena/.style={draw=fgLine, line width=0.5pt, fill=fgBand, rounded corners=1.5pt},
  %
  % ---------------------------------------------------------------- emphasis
  % Two families of the same shapes. The f family is R2's thread, in the accent
  % colour at full weight. The b family is everything else, in a light grey at
  % reduced weight. Shape, fill and outline carry the same meaning in both, so
  % the emphasis never changes what a marker means.
  %
  % robots. An open square is available to the dispatcher this epoch, a filled
  % square is working and cannot be assigned.
  fRfree/.style={rectangle, draw=fgAccent, line width=1.0pt, fill=white,
    minimum size=2.4mm, inner sep=0pt},
  fRbusy/.style={rectangle, draw=fgAccent, line width=0.7pt, fill=fgAccent,
    minimum size=2.4mm, inner sep=0pt},
  bRfree/.style={rectangle, draw=fgLine, line width=0.5pt, fill=white,
    minimum size=2.1mm, inner sep=0pt},
  bRbusy/.style={rectangle, draw=fgLine, line width=0.5pt, fill=fgLine,
    minimum size=2.1mm, inner sep=0pt},
  %
  % tasks. Circular throughout. The three states are set apart by rendering
  % rather than by stroke weight alone, so they stay separable at 1.6 mm.
  %   not yet released  one ring,  drawn at reduced opacity, ghosted
  %   pending           one ring,  full weight
  %   being carried     two rings, full weight, largest
  % Fill carries pickup against drop-off, tinted against open. The tint is kept
  % well below the outline colour so the outline stays readable on a pickup.
  % Nothing here is dashed.
  fRelP/.style={circle, draw=fgAccent, line width=1.0pt, fill=fgAccent!45,
    minimum size=1.6mm, inner sep=0pt, opacity=0.45},
  fRelD/.style={circle, draw=fgAccent, line width=1.0pt, fill=white,
    minimum size=1.6mm, inner sep=0pt, opacity=0.45},
  fPendP/.style={circle, draw=fgAccent, line width=1.0pt, fill=fgAccent!45,
    minimum size=1.7mm, inner sep=0pt},
  fPendD/.style={circle, draw=fgAccent, line width=1.0pt, fill=white,
    minimum size=1.7mm, inner sep=0pt},
  fCarryP/.style={circle, draw=fgAccent, line width=0.75pt, double=white,
    double distance=0.5pt, fill=fgAccent!45, minimum size=2.1mm, inner sep=0pt},
  fCarryD/.style={circle, draw=fgAccent, line width=0.75pt, double=white,
    double distance=0.5pt, fill=white, minimum size=2.1mm, inner sep=0pt},
  bRelP/.style={circle, draw=fgLine, line width=0.8pt, fill=fgLine!35,
    minimum size=1.5mm, inner sep=0pt, opacity=0.45},
  bRelD/.style={circle, draw=fgLine, line width=0.8pt, fill=white,
    minimum size=1.5mm, inner sep=0pt, opacity=0.45},
  bPendP/.style={circle, draw=fgLine, line width=0.8pt, fill=fgLine!35,
    minimum size=1.6mm, inner sep=0pt},
  bPendD/.style={circle, draw=fgLine, line width=0.8pt, fill=white,
    minimum size=1.6mm, inner sep=0pt},
  bCarryP/.style={circle, draw=fgLine, line width=0.5pt, double=white,
    double distance=0.4pt, fill=fgLine!35, minimum size=1.9mm, inner sep=0pt},
  bCarryD/.style={circle, draw=fgLine, line width=0.5pt, double=white,
    double distance=0.4pt, fill=white, minimum size=1.9mm, inner sep=0pt},
  %
  % the solid line pairing a pickup with its drop-off, ghosted to match while
  % the task it belongs to has not yet been released
  flink/.style={draw=fgAccent, line width=0.5pt},
  blink/.style={draw=fgLine, line width=0.3pt},
  fghost/.style={draw=fgAccent, line width=0.5pt, opacity=0.45},
  bghost/.style={draw=fgLine, line width=0.3pt, opacity=0.45},
  % the carry leg, the only dashed mark in the figure. It replaces the pairing
  % line for as long as a robot is on it.
  fcarry/.style={draw=fgAccent, line width=0.7pt,
    dash pattern=on 2.4pt off 1.5pt},
  bcarry/.style={draw=fgLine, line width=0.4pt,
    dash pattern=on 2.4pt off 1.5pt},
  % a commitment made this epoch, robot to pickup
  commit/.style={draw=fgAccent, line width=1.0pt},
  bcommit/.style={draw=fgLine, line width=0.5pt},
  % the leaders that name a pickup, a drop-off and an available robot once, in
  % the first panel, in place of legend entries for them
  leader/.style={draw=fgMid, line width=0.3pt},
  % bars
  fbar/.style={draw=fgAccent, line width=0.8pt, fill=fgAccent},
  bbar/.style={draw=fgLine, line width=0.4pt, fill=fgFill},
  guide/.style={draw=fgLine, line width=0.4pt,
    dash pattern=on 0.1pt off 1.8pt, line cap=round},
  span/.style={draw=fgAccent, line width=0.9pt},
  % text
  fTag/.style={text=fgAccent, inner sep=1.5pt, font=\figann},
  bTag/.style={text=fgMid, inner sep=1.5pt, font=\figann},
  % the single annotation set in black. It marks the point of the middle panel
  % without joining R2's accent thread, which stays reserved for R2.
  dTag/.style={text=fgDark, inner sep=1.5pt, font=\figann},
  fLab/.style={text=fgAccent, inner sep=1.5pt, font=\figann},
  bLab/.style={text=fgMid, inner sep=1.5pt, font=\figann},
  head/.style={text=fgDark, inner sep=0pt, font=\figlab}
]

\path[use as bounding box] (0,-7.55) rectangle (14.48,0.35);

% =================================================================== panels
% Panel origins. Three arenas of 4.50 cm separated by 0.49 cm, so the row is
% exactly the 14.48 cm text width. Every panel uses the same local grid of
% 10 by 11 units at 0.45 cm per unit, so a position means the same thing in all
% three and the robots can be tracked across the row.
% Four tasks. R1 takes task 1 at the first epoch, R2 takes task 2 at the first
% epoch and task 3 at the third, and R3 takes task 4 at the third. At the
% second epoch R3 is free and task 3 is pending, and R3 is left uncommitted
% anyway, which is a move a greedy assignment cannot make. R2's own position
% carries the episode.
% It starts away from the work, then sits on task 2's carry leg while it
% executes, then stands at task 2's drop-off, free again. Task 2 is delivered
% by the third panel and so is no longer drawn there.

% ------------------------------------------------------------------ epoch 1
\begin{scope}[shift={(0.00,0)}]
  \node[anchor=south west, head] at (0,0.06) {Epoch 1};
  \path[arena] (0,-4.95) rectangle (4.50,0);
  \begin{scope}[shift={(0,-4.95)}, scale=0.45]
    % pairing lines
    \draw[blink] (3.7,9.5) -- (8.7,8.4);
    \draw[flink] (2.4,4.8) -- (4.8,5.9);
    \draw[bghost] (6.6,1.6) -- (9.1,3.2);
    \draw[fghost] (6.3,6.8) -- (8.6,7.3);
    % task 1, pending, R1's
    \node[bPendP] at (3.7,9.5) {}; \node[bPendD] at (8.7,8.4) {};
    \node[bTag, anchor=south] at (3.7,9.72) {1};
    % task 2, pending, R2's
    \node[fPendP] at (2.4,4.8) {}; \node[fPendD] at (4.8,5.9) {};
    \node[fTag, anchor=east] at (2.12,4.8) {2};
    % task 4, not yet released, R3's later
    \node[bRelP] at (6.6,1.6) {}; \node[bRelD] at (9.1,3.2) {};
    \node[bTag, anchor=north] at (6.6,1.36) {4};
    % task 3, not yet released, R2's later
    \node[fRelP] at (6.3,6.8) {}; \node[fRelD] at (8.6,7.3) {};
    \node[fTag, anchor=south] at (6.3,7.04) {3};
    % commitments
    \draw[bcommit,->] (1.856,9.008) -- (3.371,9.412);
    \draw[commit,->] (1.703,3.772) -- (2.220,4.535);
    % the three marks that are named here rather than in the legend
    \draw[leader] (2.4,5.60) -- (2.4,5.02);
    \node[bTag, anchor=south] at (2.4,5.60) {pickup};
    \draw[leader] (4.8,5.30) -- (4.8,5.71);
    \node[bTag, anchor=north] at (4.8,5.30) {drop-off};
    \draw[leader] (3.6,1.05) -- (3.6,1.37);
    \node[bTag, anchor=north] at (3.6,1.05) {available};
    % robots, all three available and none of them has begun. R1 and R2 are set
    % in far enough that their labels clear the left edge of the bounding box.
    \node[bRfree] at (1.45,8.9) {}; \node[bLab, anchor=east] at (1.15,8.9) {R1};
    \node[fRfree] at (1.45,3.4) {}; \node[fLab, anchor=east] at (1.13,3.4) {R2};
    \node[bRfree] at (3.6,1.6) {}; \node[bLab, anchor=east] at (3.30,1.6) {R3};
  \end{scope}
\end{scope}

% ------------------------------------------------------------------ epoch 2
\begin{scope}[shift={(4.99,0)}]
  \node[anchor=south west, head] at (0,0.06) {Epoch 2};
  \path[arena] (0,-4.95) rectangle (4.50,0);
  \begin{scope}[shift={(0,-4.95)}, scale=0.45]
    % carry legs, the only dashed marks. R1 and R2 are on theirs.
    \draw[bcarry] (3.7,9.5) -- (8.7,8.4);
    \draw[fcarry] (2.4,4.8) -- (4.8,5.9);
    % pairing lines
    \draw[blink] (6.6,1.6) -- (9.1,3.2);
    \draw[flink] (6.3,6.8) -- (8.6,7.3);
    % task 1, being carried
    \node[bCarryP] at (3.7,9.5) {}; \node[bCarryD] at (8.7,8.4) {};
    % task 2, being carried, R2's
    \node[fCarryP] at (2.4,4.8) {}; \node[fCarryD] at (4.8,5.9) {};
    \node[fTag, anchor=east] at (2.12,4.8) {2};
    % task 4, released this epoch and left waiting
    \node[bPendP] at (6.6,1.6) {}; \node[bPendD] at (9.1,3.2) {};
    % task 3, released this epoch and left waiting, R2's later
    \node[fPendP] at (6.3,6.8) {}; \node[fPendD] at (8.6,7.3) {};
    \node[fTag, anchor=south] at (6.3,7.04) {3};
    \node[fTag, anchor=north] at (6.3,6.56) {waits};
    % no commitment is made this epoch
    % robots. R1 and R2 sit on their carry legs, so they are mid-execution and
    % cannot be assigned. R3 is at its start position and available, and is left
    % uncommitted even though task 3 is pending. Nothing here claims that
    % holding R3 back is optimal or that it is what any trained policy does.
    % It is drawn only to show that the decision is available to a dispatcher
    % and is not available to a greedy assignment.
    \node[bRbusy] at (5.7,9.06) {}; \node[bLab, anchor=south] at (5.7,9.36) {R1};
    \node[fRbusy] at (3.48,5.295) {};
    \node[fLab, anchor=north] at (3.48,5.02) {R2};
    \node[bRfree] at (3.6,1.6) {}; \node[bLab, anchor=east] at (3.30,1.6) {R3};
    \draw[leader] (3.6,1.05) -- (3.6,1.37);
    \node[dTag, anchor=north] at (3.6,1.05) {held back};
  \end{scope}
\end{scope}

% ------------------------------------------------------------------ epoch 3
\begin{scope}[shift={(9.98,0)}]
  \node[anchor=south west, head] at (0,0.06) {Epoch 3};
  \path[arena] (0,-4.95) rectangle (4.50,0);
  \begin{scope}[shift={(0,-4.95)}, scale=0.45]
    % carry legs
    \draw[bcarry] (3.7,9.5) -- (8.7,8.4);
    % pairing lines
    \draw[blink] (6.6,1.6) -- (9.1,3.2);
    \draw[flink] (6.3,6.8) -- (8.6,7.3);
    % task 1, still being carried
    \node[bCarryP] at (3.7,9.5) {}; \node[bCarryD] at (8.7,8.4) {};
    % task 2 has been delivered and is no longer drawn
    % task 4, the task R3 was held back from, taken this epoch by R3
    \node[bPendP] at (6.6,1.6) {}; \node[bPendD] at (9.1,3.2) {};
    % task 3, the waiting task, taken this epoch by R2
    \node[fPendP] at (6.3,6.8) {}; \node[fPendD] at (8.6,7.3) {};
    \node[fTag, anchor=south] at (6.3,7.04) {3};
    % commitments. R2 has just delivered, and R3 was free through the last
    % epoch and takes the task it was held back from.
    \draw[commit,->] (5.170,6.122) -- (6.035,6.641);
    \draw[bcommit,->] (4.02,1.6) -- (6.26,1.6);
    % robots. R1 is on its carry leg. R2 stands at task 2's drop-off, which is
    % where delivering it left it, and is free again. R3 has not moved.
    \node[bRbusy] at (7.7,8.62) {}; \node[bLab, anchor=south] at (7.7,8.92) {R1};
    \node[fRfree] at (4.8,5.9) {}; \node[fLab, anchor=east] at (4.48,5.9) {R2};
    \node[bRfree] at (3.6,1.6) {}; \node[bLab, anchor=east] at (3.30,1.6) {R3};
  \end{scope}
\end{scope}

% ===================================================================== bars
% Cumulative busy time after the three epochs. R2 has finished one task and is
% starting a second, so it has done the most. R1 is still on the single long
% task it took at the first epoch, and R3 stood idle through the first two
% epochs and has done the least. The bars were scaled down by a common factor
% to clear the legend column on the right, so their relative lengths are
% unchanged. The scale is 1.03 cm per unit of busy time.
\node[anchor=south west, head] at (0,-5.40) {Cumulative busy time};

\path[bbar] (0.85,-5.94) rectangle (8.53,-5.70);
\node[bLab, anchor=east] at (0.75,-5.82) {R1};

\path[fbar] (0.85,-6.36) rectangle (10.15,-6.12);
\node[fLab, anchor=east] at (0.75,-6.24) {R2};

\path[bbar] (0.85,-6.78) rectangle (5.75,-6.54);
\node[bLab, anchor=east] at (0.75,-6.66) {R3};

% the balance term, read off as the distance between the longest and the
% shortest bar. R2 is the longest, so it is one end of the span.
\draw[guide] (10.15,-6.36) -- (10.15,-7.14);
\draw[guide] (5.75,-6.78) -- (5.75,-7.14);
\draw[span,<->] (5.75,-7.14) -- (10.15,-7.14);
\node[anchor=north, inner sep=1.5pt, text=fgAccent] at (7.95,-7.20) {balance};

% =================================================================== legend
% A column of three set beside the bars rather than under them, centred on the
% middle bar so the two blocks read as one band. The three task states and
% nothing else. Pickup, drop-off and an available robot are named directly in
% the first panel instead. The swatches use the accent family, which is the
% unfaded rendering of every mark in the figure. Only the being-carried swatch
% is dashed. The column starts at 11.30, clear of the longest bar at 10.15 and
% of the balance span and its label, which sit lower and further left.
\draw[fghost] (11.30,-5.74) -- (11.90,-5.74);
\node[fRelP] at (11.30,-5.74) {}; \node[fRelD] at (11.90,-5.74) {};
\node[anchor=west, bTag] at (12.05,-5.74) {not yet released};

\draw[flink] (11.30,-6.24) -- (11.90,-6.24);
\node[fPendP] at (11.30,-6.24) {}; \node[fPendD] at (11.90,-6.24) {};
\node[anchor=west, bTag] at (12.05,-6.24) {pending};

\draw[fcarry] (11.30,-6.74) -- (11.90,-6.74);
\node[fCarryP] at (11.30,-6.74) {}; \node[fCarryD] at (11.90,-6.74) {};
\node[anchor=west, bTag] at (12.05,-6.74) {being carried};

\end{tikzpicture}%
 and do not rescale.
%
\providecommand{\figlab}{\fontsize{8}{9.4}\selectfont\hyphenpenalty=10000\relax}%
\providecommand{\figann}{\fontsize{7}{8.2}\selectfont}%
%
\providecolor{fgAccent}{RGB}{20,86,150}%
\providecolor{fgAccentFill}{RGB}{224,235,246}%
\providecolor{fgDark}{RGB}{45,45,45}%
\providecolor{fgMid}{RGB}{120,120,120}%
\providecolor{fgLine}{RGB}{160,160,160}%
\providecolor{fgFill}{RGB}{243,243,243}%
\providecolor{fgBand}{RGB}{247,247,247}%
%
\begin{tikzpicture}[
  x=1cm, y=1cm,
  >=stealth',
  line join=round,
  every node/.style={font=\figann},
  % arena frame
  arena/.style={draw=fgLine, line width=0.5pt, fill=fgBand, rounded corners=1.5pt},
  %
  % ---------------------------------------------------------------- emphasis
  % Two families of the same shapes. The f family is R2's thread, in the accent
  % colour at full weight. The b family is everything else, in a light grey at
  % reduced weight. Shape, fill and outline carry the same meaning in both, so
  % the emphasis never changes what a marker means.
  %
  % robots. An open square is available to the dispatcher this epoch, a filled
  % square is working and cannot be assigned.
  fRfree/.style={rectangle, draw=fgAccent, line width=1.0pt, fill=white,
    minimum size=2.4mm, inner sep=0pt},
  fRbusy/.style={rectangle, draw=fgAccent, line width=0.7pt, fill=fgAccent,
    minimum size=2.4mm, inner sep=0pt},
  bRfree/.style={rectangle, draw=fgLine, line width=0.5pt, fill=white,
    minimum size=2.1mm, inner sep=0pt},
  bRbusy/.style={rectangle, draw=fgLine, line width=0.5pt, fill=fgLine,
    minimum size=2.1mm, inner sep=0pt},
  %
  % tasks. Circular throughout. The three states are set apart by rendering
  % rather than by stroke weight alone, so they stay separable at 1.6 mm.
  %   not yet released  one ring,  drawn at reduced opacity, ghosted
  %   pending           one ring,  full weight
  %   being carried     two rings, full weight, largest
  % Fill carries pickup against drop-off, tinted against open. The tint is kept
  % well below the outline colour so the outline stays readable on a pickup.
  % Nothing here is dashed.
  fRelP/.style={circle, draw=fgAccent, line width=1.0pt, fill=fgAccent!45,
    minimum size=1.6mm, inner sep=0pt, opacity=0.45},
  fRelD/.style={circle, draw=fgAccent, line width=1.0pt, fill=white,
    minimum size=1.6mm, inner sep=0pt, opacity=0.45},
  fPendP/.style={circle, draw=fgAccent, line width=1.0pt, fill=fgAccent!45,
    minimum size=1.7mm, inner sep=0pt},
  fPendD/.style={circle, draw=fgAccent, line width=1.0pt, fill=white,
    minimum size=1.7mm, inner sep=0pt},
  fCarryP/.style={circle, draw=fgAccent, line width=0.75pt, double=white,
    double distance=0.5pt, fill=fgAccent!45, minimum size=2.1mm, inner sep=0pt},
  fCarryD/.style={circle, draw=fgAccent, line width=0.75pt, double=white,
    double distance=0.5pt, fill=white, minimum size=2.1mm, inner sep=0pt},
  bRelP/.style={circle, draw=fgLine, line width=0.8pt, fill=fgLine!35,
    minimum size=1.5mm, inner sep=0pt, opacity=0.45},
  bRelD/.style={circle, draw=fgLine, line width=0.8pt, fill=white,
    minimum size=1.5mm, inner sep=0pt, opacity=0.45},
  bPendP/.style={circle, draw=fgLine, line width=0.8pt, fill=fgLine!35,
    minimum size=1.6mm, inner sep=0pt},
  bPendD/.style={circle, draw=fgLine, line width=0.8pt, fill=white,
    minimum size=1.6mm, inner sep=0pt},
  bCarryP/.style={circle, draw=fgLine, line width=0.5pt, double=white,
    double distance=0.4pt, fill=fgLine!35, minimum size=1.9mm, inner sep=0pt},
  bCarryD/.style={circle, draw=fgLine, line width=0.5pt, double=white,
    double distance=0.4pt, fill=white, minimum size=1.9mm, inner sep=0pt},
  %
  % the solid line pairing a pickup with its drop-off, ghosted to match while
  % the task it belongs to has not yet been released
  flink/.style={draw=fgAccent, line width=0.5pt},
  blink/.style={draw=fgLine, line width=0.3pt},
  fghost/.style={draw=fgAccent, line width=0.5pt, opacity=0.45},
  bghost/.style={draw=fgLine, line width=0.3pt, opacity=0.45},
  % the carry leg, the only dashed mark in the figure. It replaces the pairing
  % line for as long as a robot is on it.
  fcarry/.style={draw=fgAccent, line width=0.7pt,
    dash pattern=on 2.4pt off 1.5pt},
  bcarry/.style={draw=fgLine, line width=0.4pt,
    dash pattern=on 2.4pt off 1.5pt},
  % a commitment made this epoch, robot to pickup
  commit/.style={draw=fgAccent, line width=1.0pt},
  bcommit/.style={draw=fgLine, line width=0.5pt},
  % the leaders that name a pickup, a drop-off and an available robot once, in
  % the first panel, in place of legend entries for them
  leader/.style={draw=fgMid, line width=0.3pt},
  % bars
  fbar/.style={draw=fgAccent, line width=0.8pt, fill=fgAccent},
  bbar/.style={draw=fgLine, line width=0.4pt, fill=fgFill},
  guide/.style={draw=fgLine, line width=0.4pt,
    dash pattern=on 0.1pt off 1.8pt, line cap=round},
  span/.style={draw=fgAccent, line width=0.9pt},
  % text
  fTag/.style={text=fgAccent, inner sep=1.5pt, font=\figann},
  bTag/.style={text=fgMid, inner sep=1.5pt, font=\figann},
  % the single annotation set in black. It marks the point of the middle panel
  % without joining R2's accent thread, which stays reserved for R2.
  dTag/.style={text=fgDark, inner sep=1.5pt, font=\figann},
  fLab/.style={text=fgAccent, inner sep=1.5pt, font=\figann},
  bLab/.style={text=fgMid, inner sep=1.5pt, font=\figann},
  head/.style={text=fgDark, inner sep=0pt, font=\figlab}
]

\path[use as bounding box] (0,-7.55) rectangle (14.48,0.35);

% =================================================================== panels
% Panel origins. Three arenas of 4.50 cm separated by 0.49 cm, so the row is
% exactly the 14.48 cm text width. Every panel uses the same local grid of
% 10 by 11 units at 0.45 cm per unit, so a position means the same thing in all
% three and the robots can be tracked across the row.
% Four tasks. R1 takes task 1 at the first epoch, R2 takes task 2 at the first
% epoch and task 3 at the third, and R3 takes task 4 at the third. At the
% second epoch R3 is free and task 3 is pending, and R3 is left uncommitted
% anyway, which is a move a greedy assignment cannot make. R2's own position
% carries the episode.
% It starts away from the work, then sits on task 2's carry leg while it
% executes, then stands at task 2's drop-off, free again. Task 2 is delivered
% by the third panel and so is no longer drawn there.

% ------------------------------------------------------------------ epoch 1
\begin{scope}[shift={(0.00,0)}]
  \node[anchor=south west, head] at (0,0.06) {Epoch 1};
  \path[arena] (0,-4.95) rectangle (4.50,0);
  \begin{scope}[shift={(0,-4.95)}, scale=0.45]
    % pairing lines
    \draw[blink] (3.7,9.5) -- (8.7,8.4);
    \draw[flink] (2.4,4.8) -- (4.8,5.9);
    \draw[bghost] (6.6,1.6) -- (9.1,3.2);
    \draw[fghost] (6.3,6.8) -- (8.6,7.3);
    % task 1, pending, R1's
    \node[bPendP] at (3.7,9.5) {}; \node[bPendD] at (8.7,8.4) {};
    \node[bTag, anchor=south] at (3.7,9.72) {1};
    % task 2, pending, R2's
    \node[fPendP] at (2.4,4.8) {}; \node[fPendD] at (4.8,5.9) {};
    \node[fTag, anchor=east] at (2.12,4.8) {2};
    % task 4, not yet released, R3's later
    \node[bRelP] at (6.6,1.6) {}; \node[bRelD] at (9.1,3.2) {};
    \node[bTag, anchor=north] at (6.6,1.36) {4};
    % task 3, not yet released, R2's later
    \node[fRelP] at (6.3,6.8) {}; \node[fRelD] at (8.6,7.3) {};
    \node[fTag, anchor=south] at (6.3,7.04) {3};
    % commitments
    \draw[bcommit,->] (1.856,9.008) -- (3.371,9.412);
    \draw[commit,->] (1.703,3.772) -- (2.220,4.535);
    % the three marks that are named here rather than in the legend
    \draw[leader] (2.4,5.60) -- (2.4,5.02);
    \node[bTag, anchor=south] at (2.4,5.60) {pickup};
    \draw[leader] (4.8,5.30) -- (4.8,5.71);
    \node[bTag, anchor=north] at (4.8,5.30) {drop-off};
    \draw[leader] (3.6,1.05) -- (3.6,1.37);
    \node[bTag, anchor=north] at (3.6,1.05) {available};
    % robots, all three available and none of them has begun. R1 and R2 are set
    % in far enough that their labels clear the left edge of the bounding box.
    \node[bRfree] at (1.45,8.9) {}; \node[bLab, anchor=east] at (1.15,8.9) {R1};
    \node[fRfree] at (1.45,3.4) {}; \node[fLab, anchor=east] at (1.13,3.4) {R2};
    \node[bRfree] at (3.6,1.6) {}; \node[bLab, anchor=east] at (3.30,1.6) {R3};
  \end{scope}
\end{scope}

% ------------------------------------------------------------------ epoch 2
\begin{scope}[shift={(4.99,0)}]
  \node[anchor=south west, head] at (0,0.06) {Epoch 2};
  \path[arena] (0,-4.95) rectangle (4.50,0);
  \begin{scope}[shift={(0,-4.95)}, scale=0.45]
    % carry legs, the only dashed marks. R1 and R2 are on theirs.
    \draw[bcarry] (3.7,9.5) -- (8.7,8.4);
    \draw[fcarry] (2.4,4.8) -- (4.8,5.9);
    % pairing lines
    \draw[blink] (6.6,1.6) -- (9.1,3.2);
    \draw[flink] (6.3,6.8) -- (8.6,7.3);
    % task 1, being carried
    \node[bCarryP] at (3.7,9.5) {}; \node[bCarryD] at (8.7,8.4) {};
    % task 2, being carried, R2's
    \node[fCarryP] at (2.4,4.8) {}; \node[fCarryD] at (4.8,5.9) {};
    \node[fTag, anchor=east] at (2.12,4.8) {2};
    % task 4, released this epoch and left waiting
    \node[bPendP] at (6.6,1.6) {}; \node[bPendD] at (9.1,3.2) {};
    % task 3, released this epoch and left waiting, R2's later
    \node[fPendP] at (6.3,6.8) {}; \node[fPendD] at (8.6,7.3) {};
    \node[fTag, anchor=south] at (6.3,7.04) {3};
    \node[fTag, anchor=north] at (6.3,6.56) {waits};
    % no commitment is made this epoch
    % robots. R1 and R2 sit on their carry legs, so they are mid-execution and
    % cannot be assigned. R3 is at its start position and available, and is left
    % uncommitted even though task 3 is pending. Nothing here claims that
    % holding R3 back is optimal or that it is what any trained policy does.
    % It is drawn only to show that the decision is available to a dispatcher
    % and is not available to a greedy assignment.
    \node[bRbusy] at (5.7,9.06) {}; \node[bLab, anchor=south] at (5.7,9.36) {R1};
    \node[fRbusy] at (3.48,5.295) {};
    \node[fLab, anchor=north] at (3.48,5.02) {R2};
    \node[bRfree] at (3.6,1.6) {}; \node[bLab, anchor=east] at (3.30,1.6) {R3};
    \draw[leader] (3.6,1.05) -- (3.6,1.37);
    \node[dTag, anchor=north] at (3.6,1.05) {held back};
  \end{scope}
\end{scope}

% ------------------------------------------------------------------ epoch 3
\begin{scope}[shift={(9.98,0)}]
  \node[anchor=south west, head] at (0,0.06) {Epoch 3};
  \path[arena] (0,-4.95) rectangle (4.50,0);
  \begin{scope}[shift={(0,-4.95)}, scale=0.45]
    % carry legs
    \draw[bcarry] (3.7,9.5) -- (8.7,8.4);
    % pairing lines
    \draw[blink] (6.6,1.6) -- (9.1,3.2);
    \draw[flink] (6.3,6.8) -- (8.6,7.3);
    % task 1, still being carried
    \node[bCarryP] at (3.7,9.5) {}; \node[bCarryD] at (8.7,8.4) {};
    % task 2 has been delivered and is no longer drawn
    % task 4, the task R3 was held back from, taken this epoch by R3
    \node[bPendP] at (6.6,1.6) {}; \node[bPendD] at (9.1,3.2) {};
    % task 3, the waiting task, taken this epoch by R2
    \node[fPendP] at (6.3,6.8) {}; \node[fPendD] at (8.6,7.3) {};
    \node[fTag, anchor=south] at (6.3,7.04) {3};
    % commitments. R2 has just delivered, and R3 was free through the last
    % epoch and takes the task it was held back from.
    \draw[commit,->] (5.170,6.122) -- (6.035,6.641);
    \draw[bcommit,->] (4.02,1.6) -- (6.26,1.6);
    % robots. R1 is on its carry leg. R2 stands at task 2's drop-off, which is
    % where delivering it left it, and is free again. R3 has not moved.
    \node[bRbusy] at (7.7,8.62) {}; \node[bLab, anchor=south] at (7.7,8.92) {R1};
    \node[fRfree] at (4.8,5.9) {}; \node[fLab, anchor=east] at (4.48,5.9) {R2};
    \node[bRfree] at (3.6,1.6) {}; \node[bLab, anchor=east] at (3.30,1.6) {R3};
  \end{scope}
\end{scope}

% ===================================================================== bars
% Cumulative busy time after the three epochs. R2 has finished one task and is
% starting a second, so it has done the most. R1 is still on the single long
% task it took at the first epoch, and R3 stood idle through the first two
% epochs and has done the least. The bars were scaled down by a common factor
% to clear the legend column on the right, so their relative lengths are
% unchanged. The scale is 1.03 cm per unit of busy time.
\node[anchor=south west, head] at (0,-5.40) {Cumulative busy time};

\path[bbar] (0.85,-5.94) rectangle (8.53,-5.70);
\node[bLab, anchor=east] at (0.75,-5.82) {R1};

\path[fbar] (0.85,-6.36) rectangle (10.15,-6.12);
\node[fLab, anchor=east] at (0.75,-6.24) {R2};

\path[bbar] (0.85,-6.78) rectangle (5.75,-6.54);
\node[bLab, anchor=east] at (0.75,-6.66) {R3};

% the balance term, read off as the distance between the longest and the
% shortest bar. R2 is the longest, so it is one end of the span.
\draw[guide] (10.15,-6.36) -- (10.15,-7.14);
\draw[guide] (5.75,-6.78) -- (5.75,-7.14);
\draw[span,<->] (5.75,-7.14) -- (10.15,-7.14);
\node[anchor=north, inner sep=1.5pt, text=fgAccent] at (7.95,-7.20) {balance};

% =================================================================== legend
% A column of three set beside the bars rather than under them, centred on the
% middle bar so the two blocks read as one band. The three task states and
% nothing else. Pickup, drop-off and an available robot are named directly in
% the first panel instead. The swatches use the accent family, which is the
% unfaded rendering of every mark in the figure. Only the being-carried swatch
% is dashed. The column starts at 11.30, clear of the longest bar at 10.15 and
% of the balance span and its label, which sit lower and further left.
\draw[fghost] (11.30,-5.74) -- (11.90,-5.74);
\node[fRelP] at (11.30,-5.74) {}; \node[fRelD] at (11.90,-5.74) {};
\node[anchor=west, bTag] at (12.05,-5.74) {not yet released};

\draw[flink] (11.30,-6.24) -- (11.90,-6.24);
\node[fPendP] at (11.30,-6.24) {}; \node[fPendD] at (11.90,-6.24) {};
\node[anchor=west, bTag] at (12.05,-6.24) {pending};

\draw[fcarry] (11.30,-6.74) -- (11.90,-6.74);
\node[fCarryP] at (11.30,-6.74) {}; \node[fCarryD] at (11.90,-6.74) {};
\node[anchor=west, bTag] at (12.05,-6.74) {being carried};

\end{tikzpicture}%
 and do not rescale.
%
\providecommand{\figlab}{\fontsize{8}{9.4}\selectfont\hyphenpenalty=10000\relax}%
\providecommand{\figann}{\fontsize{7}{8.2}\selectfont}%
%
\providecolor{fgAccent}{RGB}{20,86,150}%
\providecolor{fgAccentFill}{RGB}{224,235,246}%
\providecolor{fgDark}{RGB}{45,45,45}%
\providecolor{fgMid}{RGB}{120,120,120}%
\providecolor{fgLine}{RGB}{160,160,160}%
\providecolor{fgFill}{RGB}{243,243,243}%
\providecolor{fgBand}{RGB}{247,247,247}%
%
\begin{tikzpicture}[
  x=1cm, y=1cm,
  >=stealth',
  line join=round,
  every node/.style={font=\figann},
  % arena frame
  arena/.style={draw=fgLine, line width=0.5pt, fill=fgBand, rounded corners=1.5pt},
  %
  % ---------------------------------------------------------------- emphasis
  % Two families of the same shapes. The f family is R2's thread, in the accent
  % colour at full weight. The b family is everything else, in a light grey at
  % reduced weight. Shape, fill and outline carry the same meaning in both, so
  % the emphasis never changes what a marker means.
  %
  % robots. An open square is available to the dispatcher this epoch, a filled
  % square is working and cannot be assigned.
  fRfree/.style={rectangle, draw=fgAccent, line width=1.0pt, fill=white,
    minimum size=2.4mm, inner sep=0pt},
  fRbusy/.style={rectangle, draw=fgAccent, line width=0.7pt, fill=fgAccent,
    minimum size=2.4mm, inner sep=0pt},
  bRfree/.style={rectangle, draw=fgLine, line width=0.5pt, fill=white,
    minimum size=2.1mm, inner sep=0pt},
  bRbusy/.style={rectangle, draw=fgLine, line width=0.5pt, fill=fgLine,
    minimum size=2.1mm, inner sep=0pt},
  %
  % tasks. Circular throughout. The three states are set apart by rendering
  % rather than by stroke weight alone, so they stay separable at 1.6 mm.
  %   not yet released  one ring,  drawn at reduced opacity, ghosted
  %   pending           one ring,  full weight
  %   being carried     two rings, full weight, largest
  % Fill carries pickup against drop-off, tinted against open. The tint is kept
  % well below the outline colour so the outline stays readable on a pickup.
  % Nothing here is dashed.
  fRelP/.style={circle, draw=fgAccent, line width=1.0pt, fill=fgAccent!45,
    minimum size=1.6mm, inner sep=0pt, opacity=0.45},
  fRelD/.style={circle, draw=fgAccent, line width=1.0pt, fill=white,
    minimum size=1.6mm, inner sep=0pt, opacity=0.45},
  fPendP/.style={circle, draw=fgAccent, line width=1.0pt, fill=fgAccent!45,
    minimum size=1.7mm, inner sep=0pt},
  fPendD/.style={circle, draw=fgAccent, line width=1.0pt, fill=white,
    minimum size=1.7mm, inner sep=0pt},
  fCarryP/.style={circle, draw=fgAccent, line width=0.75pt, double=white,
    double distance=0.5pt, fill=fgAccent!45, minimum size=2.1mm, inner sep=0pt},
  fCarryD/.style={circle, draw=fgAccent, line width=0.75pt, double=white,
    double distance=0.5pt, fill=white, minimum size=2.1mm, inner sep=0pt},
  bRelP/.style={circle, draw=fgLine, line width=0.8pt, fill=fgLine!35,
    minimum size=1.5mm, inner sep=0pt, opacity=0.45},
  bRelD/.style={circle, draw=fgLine, line width=0.8pt, fill=white,
    minimum size=1.5mm, inner sep=0pt, opacity=0.45},
  bPendP/.style={circle, draw=fgLine, line width=0.8pt, fill=fgLine!35,
    minimum size=1.6mm, inner sep=0pt},
  bPendD/.style={circle, draw=fgLine, line width=0.8pt, fill=white,
    minimum size=1.6mm, inner sep=0pt},
  bCarryP/.style={circle, draw=fgLine, line width=0.5pt, double=white,
    double distance=0.4pt, fill=fgLine!35, minimum size=1.9mm, inner sep=0pt},
  bCarryD/.style={circle, draw=fgLine, line width=0.5pt, double=white,
    double distance=0.4pt, fill=white, minimum size=1.9mm, inner sep=0pt},
  %
  % the solid line pairing a pickup with its drop-off, ghosted to match while
  % the task it belongs to has not yet been released
  flink/.style={draw=fgAccent, line width=0.5pt},
  blink/.style={draw=fgLine, line width=0.3pt},
  fghost/.style={draw=fgAccent, line width=0.5pt, opacity=0.45},
  bghost/.style={draw=fgLine, line width=0.3pt, opacity=0.45},
  % the carry leg, the only dashed mark in the figure. It replaces the pairing
  % line for as long as a robot is on it.
  fcarry/.style={draw=fgAccent, line width=0.7pt,
    dash pattern=on 2.4pt off 1.5pt},
  bcarry/.style={draw=fgLine, line width=0.4pt,
    dash pattern=on 2.4pt off 1.5pt},
  % a commitment made this epoch, robot to pickup
  commit/.style={draw=fgAccent, line width=1.0pt},
  bcommit/.style={draw=fgLine, line width=0.5pt},
  % the leaders that name a pickup, a drop-off and an available robot once, in
  % the first panel, in place of legend entries for them
  leader/.style={draw=fgMid, line width=0.3pt},
  % bars
  fbar/.style={draw=fgAccent, line width=0.8pt, fill=fgAccent},
  bbar/.style={draw=fgLine, line width=0.4pt, fill=fgFill},
  guide/.style={draw=fgLine, line width=0.4pt,
    dash pattern=on 0.1pt off 1.8pt, line cap=round},
  span/.style={draw=fgAccent, line width=0.9pt},
  % text
  fTag/.style={text=fgAccent, inner sep=1.5pt, font=\figann},
  bTag/.style={text=fgMid, inner sep=1.5pt, font=\figann},
  % the single annotation set in black. It marks the point of the middle panel
  % without joining R2's accent thread, which stays reserved for R2.
  dTag/.style={text=fgDark, inner sep=1.5pt, font=\figann},
  fLab/.style={text=fgAccent, inner sep=1.5pt, font=\figann},
  bLab/.style={text=fgMid, inner sep=1.5pt, font=\figann},
  head/.style={text=fgDark, inner sep=0pt, font=\figlab}
]

\path[use as bounding box] (0,-7.55) rectangle (14.48,0.35);

% =================================================================== panels
% Panel origins. Three arenas of 4.50 cm separated by 0.49 cm, so the row is
% exactly the 14.48 cm text width. Every panel uses the same local grid of
% 10 by 11 units at 0.45 cm per unit, so a position means the same thing in all
% three and the robots can be tracked across the row.
% Four tasks. R1 takes task 1 at the first epoch, R2 takes task 2 at the first
% epoch and task 3 at the third, and R3 takes task 4 at the third. At the
% second epoch R3 is free and task 3 is pending, and R3 is left uncommitted
% anyway, which is a move a greedy assignment cannot make. R2's own position
% carries the episode.
% It starts away from the work, then sits on task 2's carry leg while it
% executes, then stands at task 2's drop-off, free again. Task 2 is delivered
% by the third panel and so is no longer drawn there.

% ------------------------------------------------------------------ epoch 1
\begin{scope}[shift={(0.00,0)}]
  \node[anchor=south west, head] at (0,0.06) {Epoch 1};
  \path[arena] (0,-4.95) rectangle (4.50,0);
  \begin{scope}[shift={(0,-4.95)}, scale=0.45]
    % pairing lines
    \draw[blink] (3.7,9.5) -- (8.7,8.4);
    \draw[flink] (2.4,4.8) -- (4.8,5.9);
    \draw[bghost] (6.6,1.6) -- (9.1,3.2);
    \draw[fghost] (6.3,6.8) -- (8.6,7.3);
    % task 1, pending, R1's
    \node[bPendP] at (3.7,9.5) {}; \node[bPendD] at (8.7,8.4) {};
    \node[bTag, anchor=south] at (3.7,9.72) {1};
    % task 2, pending, R2's
    \node[fPendP] at (2.4,4.8) {}; \node[fPendD] at (4.8,5.9) {};
    \node[fTag, anchor=east] at (2.12,4.8) {2};
    % task 4, not yet released, R3's later
    \node[bRelP] at (6.6,1.6) {}; \node[bRelD] at (9.1,3.2) {};
    \node[bTag, anchor=north] at (6.6,1.36) {4};
    % task 3, not yet released, R2's later
    \node[fRelP] at (6.3,6.8) {}; \node[fRelD] at (8.6,7.3) {};
    \node[fTag, anchor=south] at (6.3,7.04) {3};
    % commitments
    \draw[bcommit,->] (1.856,9.008) -- (3.371,9.412);
    \draw[commit,->] (1.703,3.772) -- (2.220,4.535);
    % the three marks that are named here rather than in the legend
    \draw[leader] (2.4,5.60) -- (2.4,5.02);
    \node[bTag, anchor=south] at (2.4,5.60) {pickup};
    \draw[leader] (4.8,5.30) -- (4.8,5.71);
    \node[bTag, anchor=north] at (4.8,5.30) {drop-off};
    \draw[leader] (3.6,1.05) -- (3.6,1.37);
    \node[bTag, anchor=north] at (3.6,1.05) {available};
    % robots, all three available and none of them has begun. R1 and R2 are set
    % in far enough that their labels clear the left edge of the bounding box.
    \node[bRfree] at (1.45,8.9) {}; \node[bLab, anchor=east] at (1.15,8.9) {R1};
    \node[fRfree] at (1.45,3.4) {}; \node[fLab, anchor=east] at (1.13,3.4) {R2};
    \node[bRfree] at (3.6,1.6) {}; \node[bLab, anchor=east] at (3.30,1.6) {R3};
  \end{scope}
\end{scope}

% ------------------------------------------------------------------ epoch 2
\begin{scope}[shift={(4.99,0)}]
  \node[anchor=south west, head] at (0,0.06) {Epoch 2};
  \path[arena] (0,-4.95) rectangle (4.50,0);
  \begin{scope}[shift={(0,-4.95)}, scale=0.45]
    % carry legs, the only dashed marks. R1 and R2 are on theirs.
    \draw[bcarry] (3.7,9.5) -- (8.7,8.4);
    \draw[fcarry] (2.4,4.8) -- (4.8,5.9);
    % pairing lines
    \draw[blink] (6.6,1.6) -- (9.1,3.2);
    \draw[flink] (6.3,6.8) -- (8.6,7.3);
    % task 1, being carried
    \node[bCarryP] at (3.7,9.5) {}; \node[bCarryD] at (8.7,8.4) {};
    % task 2, being carried, R2's
    \node[fCarryP] at (2.4,4.8) {}; \node[fCarryD] at (4.8,5.9) {};
    \node[fTag, anchor=east] at (2.12,4.8) {2};
    % task 4, released this epoch and left waiting
    \node[bPendP] at (6.6,1.6) {}; \node[bPendD] at (9.1,3.2) {};
    % task 3, released this epoch and left waiting, R2's later
    \node[fPendP] at (6.3,6.8) {}; \node[fPendD] at (8.6,7.3) {};
    \node[fTag, anchor=south] at (6.3,7.04) {3};
    \node[fTag, anchor=north] at (6.3,6.56) {waits};
    % no commitment is made this epoch
    % robots. R1 and R2 sit on their carry legs, so they are mid-execution and
    % cannot be assigned. R3 is at its start position and available, and is left
    % uncommitted even though task 3 is pending. Nothing here claims that
    % holding R3 back is optimal or that it is what any trained policy does.
    % It is drawn only to show that the decision is available to a dispatcher
    % and is not available to a greedy assignment.
    \node[bRbusy] at (5.7,9.06) {}; \node[bLab, anchor=south] at (5.7,9.36) {R1};
    \node[fRbusy] at (3.48,5.295) {};
    \node[fLab, anchor=north] at (3.48,5.02) {R2};
    \node[bRfree] at (3.6,1.6) {}; \node[bLab, anchor=east] at (3.30,1.6) {R3};
    \draw[leader] (3.6,1.05) -- (3.6,1.37);
    \node[dTag, anchor=north] at (3.6,1.05) {held back};
  \end{scope}
\end{scope}

% ------------------------------------------------------------------ epoch 3
\begin{scope}[shift={(9.98,0)}]
  \node[anchor=south west, head] at (0,0.06) {Epoch 3};
  \path[arena] (0,-4.95) rectangle (4.50,0);
  \begin{scope}[shift={(0,-4.95)}, scale=0.45]
    % carry legs
    \draw[bcarry] (3.7,9.5) -- (8.7,8.4);
    % pairing lines
    \draw[blink] (6.6,1.6) -- (9.1,3.2);
    \draw[flink] (6.3,6.8) -- (8.6,7.3);
    % task 1, still being carried
    \node[bCarryP] at (3.7,9.5) {}; \node[bCarryD] at (8.7,8.4) {};
    % task 2 has been delivered and is no longer drawn
    % task 4, the task R3 was held back from, taken this epoch by R3
    \node[bPendP] at (6.6,1.6) {}; \node[bPendD] at (9.1,3.2) {};
    % task 3, the waiting task, taken this epoch by R2
    \node[fPendP] at (6.3,6.8) {}; \node[fPendD] at (8.6,7.3) {};
    \node[fTag, anchor=south] at (6.3,7.04) {3};
    % commitments. R2 has just delivered, and R3 was free through the last
    % epoch and takes the task it was held back from.
    \draw[commit,->] (5.170,6.122) -- (6.035,6.641);
    \draw[bcommit,->] (4.02,1.6) -- (6.26,1.6);
    % robots. R1 is on its carry leg. R2 stands at task 2's drop-off, which is
    % where delivering it left it, and is free again. R3 has not moved.
    \node[bRbusy] at (7.7,8.62) {}; \node[bLab, anchor=south] at (7.7,8.92) {R1};
    \node[fRfree] at (4.8,5.9) {}; \node[fLab, anchor=east] at (4.48,5.9) {R2};
    \node[bRfree] at (3.6,1.6) {}; \node[bLab, anchor=east] at (3.30,1.6) {R3};
  \end{scope}
\end{scope}

% ===================================================================== bars
% Cumulative busy time after the three epochs. R2 has finished one task and is
% starting a second, so it has done the most. R1 is still on the single long
% task it took at the first epoch, and R3 stood idle through the first two
% epochs and has done the least. The bars were scaled down by a common factor
% to clear the legend column on the right, so their relative lengths are
% unchanged. The scale is 1.03 cm per unit of busy time.
\node[anchor=south west, head] at (0,-5.40) {Cumulative busy time};

\path[bbar] (0.85,-5.94) rectangle (8.53,-5.70);
\node[bLab, anchor=east] at (0.75,-5.82) {R1};

\path[fbar] (0.85,-6.36) rectangle (10.15,-6.12);
\node[fLab, anchor=east] at (0.75,-6.24) {R2};

\path[bbar] (0.85,-6.78) rectangle (5.75,-6.54);
\node[bLab, anchor=east] at (0.75,-6.66) {R3};

% the balance term, read off as the distance between the longest and the
% shortest bar. R2 is the longest, so it is one end of the span.
\draw[guide] (10.15,-6.36) -- (10.15,-7.14);
\draw[guide] (5.75,-6.78) -- (5.75,-7.14);
\draw[span,<->] (5.75,-7.14) -- (10.15,-7.14);
\node[anchor=north, inner sep=1.5pt, text=fgAccent] at (7.95,-7.20) {balance};

% =================================================================== legend
% A column of three set beside the bars rather than under them, centred on the
% middle bar so the two blocks read as one band. The three task states and
% nothing else. Pickup, drop-off and an available robot are named directly in
% the first panel instead. The swatches use the accent family, which is the
% unfaded rendering of every mark in the figure. Only the being-carried swatch
% is dashed. The column starts at 11.30, clear of the longest bar at 10.15 and
% of the balance span and its label, which sit lower and further left.
\draw[fghost] (11.30,-5.74) -- (11.90,-5.74);
\node[fRelP] at (11.30,-5.74) {}; \node[fRelD] at (11.90,-5.74) {};
\node[anchor=west, bTag] at (12.05,-5.74) {not yet released};

\draw[flink] (11.30,-6.24) -- (11.90,-6.24);
\node[fPendP] at (11.30,-6.24) {}; \node[fPendD] at (11.90,-6.24) {};
\node[anchor=west, bTag] at (12.05,-6.24) {pending};

\draw[fcarry] (11.30,-6.74) -- (11.90,-6.74);
\node[fCarryP] at (11.30,-6.74) {}; \node[fCarryD] at (11.90,-6.74) {};
\node[anchor=west, bTag] at (12.05,-6.74) {being carried};

\end{tikzpicture}%
